\chapter{Installation and first run}
\label{ch:install}

\section{Requirements}

\begin{center}
\begin{tabularx}{\textwidth}{@{}lX@{}}
\toprule
Item & Requirement \\
\midrule
Python & 3.11 or newer (3.12 is used for development) \\
Operating system & Windows, Linux or macOS \\
Memory & 2\,GB free is ample \\
Runtime libraries & numpy, pandas, pymodbus 3.8.x, asyncua, duckdb, fastapi, uvicorn, httpx \\
Machine learning (\code{ml} extra) & scikit-learn, pyarrow, joblib \\
Workbench UI (\code{ui} extra) & dash, plotly \\
Tests (\code{dev} extra) & pytest \\
\bottomrule
\end{tabularx}
\end{center}

\begin{warn}[pymodbus must stay on 3.8.x]
The 3.15 line of pymodbus replaced the datastore API this project uses and no
longer offers a way to update served values from outside the server, which is
what a live plant does twenty times a second. \code{pyproject.toml} pins
\code{pymodbus>=3.8,<3.9}; do not relax it.
\end{warn}

\section{Installing}

From the directory that contains \code{pyproject.toml}:

\begin{lstlisting}[style=sh]
# a dedicated environment is recommended
conda create -n pyplant python=3.12 -y
conda activate pyplant

pip uninstall -y pyplant          # an unrelated PyPI package has this name
pip install -e ".[ml,ui,dev]"     # everything: stack, ML, workbench, tests
python -c "import pyplant.hist.db; print('ok')"
\end{lstlisting}

The \code{-e} (editable) flag puts this directory on the Python path, so an
edited file takes effect without reinstalling and \code{python -m pyplant...}
works from any directory.

\begin{note}[Windows with several Pythons]
If \code{python --version} reports something older than 3.11, call the
environment's interpreter explicitly, for example
\begin{lstlisting}[style=sh]
C:\Users\you\anaconda3\envs\pyplant\python.exe -m pyplant.workbench.app
\end{lstlisting}
\end{note}

\section{Checking the installation}

\begin{lstlisting}[style=sh]
python -m pyplant.tags          # prints the tag list and the Modbus block plan
pytest -q                       # full suite, about five minutes
pytest -q -m "not slow"         # skips tests that start real subprocesses
\end{lstlisting}

\section{Starting the workbench}

\begin{lstlisting}[style=sh]
python -m pyplant.workbench.app --workspace ./work/demo --base-port 5400
\end{lstlisting}

Then open \url{http://127.0.0.1:8050}. The options are:

\begin{center}
\begin{tabular}{@{}lll@{}}
\toprule
Option & Default & Meaning \\
\midrule
\code{--workspace} & \code{./work/demo} & directory holding config, data, models and logs \\
\code{--base-port} & 5400 & first of six consecutive ports for the stack \\
\code{--host} & 127.0.0.1 & interface the workbench page listens on \\
\code{--port} & 8050 & port of the workbench page \\
\code{--debug} & off & Dash debug mode with hot reload \\
\bottomrule
\end{tabular}
\end{center}

With \code{--base-port 5400} the stack uses the ports in
Table~\ref{tab:ports}. Several people on one machine can each run a stack by
choosing different base ports.

\begin{table}[htbp]
\centering
\begin{tabular}{@{}llll@{}}
\toprule
Port & Offset & Service & Protocol \\
\midrule
5400 & +0 & plant Modbus slave & Modbus TCP \\
5401 & +1 & plant fault-injection API & HTTP \\
5402 & +2 & PLC Modbus server (SCADA side) & Modbus TCP \\
5403 & +3 & PLC OPC UA endpoint, \code{/pyplant/} & OPC UA \\
5404 & +4 & historian query API & HTTP \\
5405 & +5 & HMI operator page and API & HTTP \\
\bottomrule
\end{tabular}
\caption{Ports allocated from the base port.}
\label{tab:ports}
\end{table}

\section{The workspace directory}

Everything a session produces lives in the workspace, and nothing it produces
lives anywhere else:

\begin{lstlisting}[style=sh]
work/demo/
  config.json          speed, dt, ports, plant, loop and data settings
  runs.jsonl           append-only log of every action and its result
  dataset.parquet      the generated, labelled dataset
  history.duckdb       the historian database of the live stack
  artifacts/           trained detectors, manifest.json, metrics.json, per-run CSVs
  jobs/                progress files of background jobs and service logs
\end{lstlisting}

Deleting the workspace is always safe: the next start recreates it with
default settings.

\section{Running the stack by hand}

The workbench starts the four services for you, but each can be started from
a terminal. This is useful for debugging and is exactly what the workbench
runs (the Runtime tab shows the command lines).

\begin{lstlisting}[style=sh]
python -m pyplant.services.plant     --port 5020 --control-port 5021
python -m pyplant.services.plc       --plant-port 5020 --port 5022 --opcua-port 5023
python -m pyplant.services.historian --endpoint opc.tcp://127.0.0.1:5023/pyplant/ \
                                     --db history.duckdb --api-port 5024
python -m pyplant.services.hmi       --endpoint opc.tcp://127.0.0.1:5023/pyplant/ \
                                     --history http://127.0.0.1:5024 --port 5025
\end{lstlisting}

Start them in this order; each depends on the one before it. All command-line
options are listed in Section~\ref{sec:cli}.
